
\section*{Abstract}

We present \textbf{AutonomousDiscoveryLoop}, the orchestration component that closes the end-to-end autonomous knowledge graph improvement loop in CEREBRUM. The loop runs \texttt{ResearchAgent.scan\_once()} on a configurable timer, processes each finding through \texttt{AutoApprover} (Phase 71), and applies a \textbf{circuit breaker} --- a sliding-window approval-rate monitor that pauses materialization when quality degrades below a threshold. A \textbf{per-cycle cap} (\texttt{max\_materializations\_per\_cycle}) prevents runaway edge insertion during anomalous discovery bursts. A \textbf{dry_run} mode allows safe production trials without writing any edges. \texttt{AutoApprover} state is checkpointed to disk after every cycle, enabling warm restart. The loop exposes a structured REST API for monitoring, configuration, and lifecycle management, making it suitable for embedding in production CEREBRUM deployments without operator intervention.

\medskip\noindent\rule{\linewidth}{0.4pt}\medskip

\section*{1. Architecture}

\subsection*{1.1 Loop Lifecycle}

\begin{lstlisting}
start() → background thread
    ↓
while running:
    scan_once() → findings
    for finding in findings[:cap]:
        decide = approver.decide(finding)
        if decide == APPROVE: approve(finding)
        if decide == REJECT:  reject(finding)
    record_cycle()
    checkpoint_approver()
    check_circuit_breaker()
    sleep(next_interval)
stop() → graceful shutdown
\end{lstlisting}

\subsection*{1.2 LoopConfig Dataclass}

\begin{lstlisting}[language=Python]
@dataclass
class LoopConfig:
    cycle_interval: float = 300.0
    max_materializations_per_cycle: int = 10
    min_approval_rate: float = 0.5
    circuit_breaker_window: int = 20
    dry_run: bool = False
    auto_rollback_on_trip: bool = False   # Phase 79
    adaptive_tuning: bool = False          # Phase 82
    adaptive_min_cap: int = 1
    adaptive_max_cap: int = 20
    adaptive_min_interval: float = 60.0
    adaptive_max_interval: float = 7200.0
    approver_checkpoint_path: Optional[str] = None
\end{lstlisting}

\subsection*{1.3 CycleRecord Dataclass}

\begin{lstlisting}[language=Python]
@dataclass
class CycleRecord:
    cycle_number: int
    started_at: float
    findings_scanned: int
    materializations: int
    approvals: int
    rejections: int
    reviews: int
    circuit_breaker_tripped: bool
    edges_rolled_back: int = 0    # Phase 79
    effective_cap: int = 0        # Phase 82
\end{lstlisting}

\medskip\noindent\rule{\linewidth}{0.4pt}\medskip

\section*{2. Circuit Breaker}

The circuit breaker computes approval rate over a sliding window of the last N decisions (approvals + rejections):

\begin{lstlisting}
approval_rate = approvals_in_window / (approvals + rejections)_in_window
\end{lstlisting}

If \texttt{approval\_rate < min\_approval\_rate}:
\begin{itemize}
  \item \texttt{circuit\_breaker\_tripped = True}
  \item Materialization pauses for the current cycle
  \item If \texttt{auto\_rollback\_on\_trip=True} (Phase 79): \texttt{ProvenanceLedger.rollback\_cycle()} is called
  \item Loop continues sleeping; next cycle attempts recovery
\end{itemize}


The circuit breaker resets when the approval rate recovers above threshold over a fresh window.

\medskip\noindent\rule{\linewidth}{0.4pt}\medskip

\section*{3. Per-Cycle Cap}

\texttt{max\_materializations\_per\_cycle} acts as a hard upper bound on approved findings per cycle. Even if \texttt{AutoApprover} would approve 100 findings in a single scan, only the first N are materialized. This prevents:
\begin{itemize}
  \item Burst materialization that overwhelms downstream systems
  \item Runaway edge accumulation during transient high-discovery phases
\end{itemize}


Combined with \texttt{adaptive\_tuning} (Phase 82), the cap is dynamically scaled per cycle based on \texttt{DiscoveryCalibrator} community weights.

\medskip\noindent\rule{\linewidth}{0.4pt}\medskip

\section*{4. Dry Run Mode}

\texttt{LoopConfig(dry\_run=True)} makes the loop execute all phases (scan, validate, decide) but skips \texttt{approve()} / \texttt{reject()} calls. The loop records what \textit{would} have been materialized in \texttt{CycleRecord}, enabling safe production trials to measure expected impact before enabling writes.

\medskip\noindent\rule{\linewidth}{0.4pt}\medskip

\section*{5. REST API}

\begin{table}[h]
\small\centering
\begin{tabular}{lll}
\toprule
Endpoint & Method & Description \\
\midrule
\texttt{/research/loop/start} & POST & Start loop (idempotent) \\
\texttt{/research/loop/stop} & POST & Graceful stop \\
\texttt{/research/loop/status} & GET & Running state, cycle history, approval rate, circuit breaker \\
\texttt{/research/loop/configure} & POST & Partial update: \texttt{cycle\_interval}, \texttt{max\_materializations\_per\_cycle}, \texttt{min\_approval\_rate}, \texttt{dry\_run}, \texttt{auto\_rollback\_on\_trip}, \texttt{adaptive\_tuning} \\
\bottomrule
\end{tabular}
\end{table}


\medskip\noindent\rule{\linewidth}{0.4pt}\medskip

\section*{6. References}

\begin{itemize}
  \item [dean2008mapreduce] Dean, J. \& Ghemawat, S. (2008). MapReduce: Simplified data processing on large clusters. \textit{CACM}, 51(1), 107--113.
\end{itemize}
